% Notes of the Montreal seminar, 1970 (cote 129, typed part): archive cote 129, site page 5 (PDF page 6).
% typescript. Draft compared: none.
% AI-transcribed by Claude Opus 5.5 (claude-opus-5-5, adaptive thinking, effort high), Message Batches API, batch mtl129_full (msgbatch_01AGnDKHwy6ATCM7ehUa6263), 2026-11-13; not yet checked by eye.
\scanpage{129/5}
Proposition (2.3). En caractéristique p, on a les relations
\[ F.V = V.F = p\ \mathrm{Id}_{\mathbb{W}_{\mathbb{F}_p}} \]
(c.f. J.-P. Serre "corps locaux" II §6 Cor. du th.7). On peut
donner la démonstration directe suivante: Posons
\[ h = (h_1,h_2,\ldots,h_n,\ldots) = (p.\mathrm{Id}_W - VF)(x_1,\ldots,x_n,\ldots). \]
On sait que $h_i \in \mathbb{Z}[x_1,x_2,\ldots,x_i]$, et il suffit de démontrer que $h_i$
est divisible par $p$ pour tout $i$ pour prouver la proposition. Or
on trouve que
\[ \phi_i(h) = p\phi_i(x) - \phi_i(VF(x)) = p^i\, x_i \]
d'où l'on déduit que $h_1 = px_1$,
\[ \phi_2(h) = h_1^p + ph_2 = p^2\, x_2 \quad \text{d'où} \quad h_2 = p(x_2-p^{p-2}\, x_1), \]
et par récurrence sur $i$, on en déduit que pour tout $i$ on a $h_i = pk_i$.

Notons en particulier que dans $\mathbb{W}_{\mathbb{F}_p}$ on a $p = (0,1,0,\ldots)$
(car $(1,0,0,\ldots)$ est l'unité de $\mathbb{W}$).

3. \uline{L'ind-schéma} $\underset{\rightarrow}{\mathbb{W}}$ \uline{et l'ind-schéma} $\underset{\rightarrow}{\mathbb{W}}'$

On définit
\[ \underset{\rightarrow}{\mathbb{W}} = \underset{n}{\underset{\rightarrow}{\lim}}\, (\mathbb{W}_n,T_n). \]
C'est un ind-schéma en groupes dont le ind-schéma sous-jacent est

\mspage{\unsure{5}}
